\documentclass[a4paper,12pt]{article}

\usepackage{iftex}
\ifPDFTeX
  \usepackage[T2A]{fontenc}
  \usepackage[utf8]{inputenc}
  \usepackage[russian]{babel}
  \usepackage{helvet}
  \usepackage{microtype}
\else
  \usepackage{fontspec}
  \setmainfont{DejaVu Sans}[AutoFakeSlant=0.18]
  \setsansfont{DejaVu Sans}[AutoFakeSlant=0.18]
\fi
\renewcommand{\familydefault}{\sfdefault}
\renewcommand{\today}{17 июля 2026 г.}
\usepackage{amsmath,amssymb,amsthm,mathtools}
\usepackage{icomma}
\usepackage[a4paper,left=18mm,right=18mm,top=18mm,bottom=20mm]{geometry}
\usepackage{tikz}
\usepackage{tkz-euclide}
\usepackage{pgfplots}
\pgfplotsset{compat=1.18}
\usepackage{tabularx,booktabs,longtable,array}
\usepackage{enumitem}
\usepackage{hyperref}
\usepackage{parskip}
\usepackage{fancyhdr}
\usepackage{setspace}
\usepackage{xcolor}

\definecolor{accent}{HTML}{7B3142}
\definecolor{darkaccent}{HTML}{4E2430}
\definecolor{textgray}{HTML}{30343B}
\definecolor{softgray}{HTML}{6E7279}

\hypersetup{colorlinks=true,linkcolor=darkaccent,urlcolor=accent}
\onehalfspacing
\setlength{\parindent}{0pt}
\setlength{\parskip}{5pt}
\setlength{\emergencystretch}{3em}
\setlist[itemize]{leftmargin=6mm,itemsep=2pt,topsep=3pt}
\setlist[enumerate]{leftmargin=8mm,itemsep=4pt,topsep=4pt}
\renewcommand{\arraystretch}{1.35}

\pagestyle{fancy}
\fancyhf{}
\lhead{\textcolor{textgray}{Показательные неравенства: уровень выше}}
\rhead{\textcolor{textgray}{17.07.26}}
\cfoot{\textcolor{softgray}{\thepage}}
\setlength{\headheight}{15pt}

\newcommand{\key}[1]{\textcolor{accent}{\textbf{#1}}}
\newcommand{\R}{\mathbb{R}}

\begin{document}
\sloppy

\begin{titlepage}
  \centering
  \vspace*{25mm}
  {\Large\textcolor{accent}{ЧЕК-ЛИСТ}\par}
  \vspace{9mm}
  {\fontsize{29}{35}\selectfont\bfseries Показательные неравенства:\\[2mm] уровень выше\par}
  \vspace{8mm}
  {\large Модуль, системы и совокупности,\\группировка и разность степеней\par}
  \vfill
  {\large Лёвин Артём Александрович\\эксклюзивно для Тимофея Васильченко\par}
  \vspace{8mm}
  {\large \today\par}
  \vfill
  \begin{tikzpicture}[remember picture,overlay]
    \draw[accent,line width=2pt]
      ([xshift=22mm,yshift=24mm]current page.south west)
      .. controls +(38mm,22mm) and +(-35mm,-10mm) ..
      ([xshift=-22mm,yshift=35mm]current page.south east);
    \draw[darkaccent,line width=.8pt]
      ([xshift=32mm,yshift=18mm]current page.south west)
      .. controls +(45mm,28mm) and +(-48mm,-14mm) ..
      ([xshift=-30mm,yshift=28mm]current page.south east);
  \end{tikzpicture}
\end{titlepage}

\section{Карта занятия}

На прошлом занятии были освоены общее основание и замена переменной. Теперь задача --- распознавать более сложную структуру: модуль в показателе, четыре слагаемых, произведение разностей степеней и деление на показательное выражение.

\begin{enumerate}[itemsep=1pt]
  \item привести обе части к одинаковому основанию или к сравнению с единицей;
  \item установить, возрастает или убывает показательная функция;
  \item снять показатель и решить полученное алгебраическое неравенство;
  \item при четырёх слагаемых проверить группировку и двойное вынесение;
  \item делить только на выражение с заранее известным знаком;
  \item проверить строгие границы, ОДЗ и запись объединения промежутков.
\end{enumerate}

\section{Основание и сравнение с единицей}

Для $a>0$, $a\ne1$ справедливо
\[
 a^{f(x)}\gtreqless a^{g(x)}
 \iff
 \begin{cases}
 f(x)\gtreqless g(x),&a>1,\\
 f(x)\lesseqgtr g(x),&0<a<1.
 \end{cases}
\]

Особенно полезно представление $1=a^0$. Перед снятием основания обязательно установите, по какую сторону от единицы находится $a$.

\begin{table}[ht]
\centering
\caption{Что происходит со знаком неравенства}
\begin{tabularx}{\linewidth}{>{\bfseries}p{31mm} X X}
\toprule
Основание & Монотонность & Действие \\
\midrule
$a>1$ & функция $a^x$ возрастает & знак между показателями сохраняется \\
$0<a<1$ & функция $a^x$ убывает & знак между показателями меняется \\
\bottomrule
\end{tabularx}
\end{table}

\subsection{Пример с модулем в показателе}

Решим
\[
\left(\frac{\pi}{3}\right)^{|x^2-x|-3}\ge1.
\]
Поскольку $\pi>3$, имеем $\pi/3>1$, поэтому знак сохраняется:
\[
|x^2-x|-3\ge0\iff |x^2-x|\ge3.
\]
Раскрываем модуль как \key{совокупность}:
\[
\left[
\begin{aligned}
x^2-x&\ge3,\\
x^2-x&\le-3.
\end{aligned}
\right.
\]
Первое неравенство даёт
\[
x^2-x-3\ge0,qquad
x\in\left(-\infty;\frac{1-\sqrt{13}}2\right]
\cup
\left[\frac{1+\sqrt{13}}2;+\infty\right).
\]
У второго неравенства $x^2-x+3\le0$ дискриминант равен $-11$, а квадратный трёхчлен положителен при всех $x$. Поэтому решений оно не имеет. Итог:
\[
\boxed{x\in\left(-\infty;\frac{1-\sqrt{13}}2\right]
\cup
\left[\frac{1+\sqrt{13}}2;+\infty\right)}.
\]

\section{Модуль: система или совокупность}

После изоляции модуля используйте равносильный переход, соответствующий знаку.
\[
|f(x)|\ge g(x)
\iff
\left[
\begin{aligned}
f(x)&\ge g(x),\\
f(x)&\le-g(x),
\end{aligned}
\right.
\]
\[
|f(x)|\le g(x)
\iff
\begin{cases}
f(x)\le g(x),\\
f(x)\ge-g(x).
\end{cases}
\]
Для строгих знаков $>$ и $<$ строгими остаются обе ветви. Совокупность объединяет решения хотя бы одной ветви; система берёт их пересечение.

\subsection{Два коротких примера}

\[
|x+1|-3x+5\ge0
\iff |x+1|\ge3x-5
\iff
\left[
\begin{aligned}
x+1&\ge3x-5,\\
x+1&\le-3x+5.
\end{aligned}
\right.
\]
Отсюда $x\le3$ или $x\le1$, то есть $\boxed{x\le3}$.

\[
|x+1|-3x+5<0
\iff |x+1|<3x-5
\iff
\begin{cases}
x+1<3x-5,\\
x+1>-3x+5.
\end{cases}
\]
Получаем $x>3$ и $x>1$, следовательно $\boxed{x>3}$.

\key{Самопроверка:} если справа от модуля получилось отрицательное число, то неравенство вида $|f(x)|\le g(x)$ решений не имеет, а $|f(x)|\ge g(x)$ выполняется автоматически. При переменном $g(x)$ это нужно учитывать через систему или совокупность, а не отдельной догадкой.

\section{Группировка и двойное вынесение}

Четыре слагаемых часто являются сигналом к группировке. Разберём пример занятия:
\[
25\cdot2^x-10^x+5^x-25>0.
\]
Так как $10^x=2^x5^x$, группируем:
\begin{align*}
25\cdot2^x-2^x5^x+5^x-25
&=2^x(25-5^x)-(25-5^x)\\
&=(25-5^x)(2^x-1).
\end{align*}

Это и есть \key{двойное вынесение}: сначала общий множитель появляется внутри каждой группы, затем одинаковая скобка выносится из всего выражения.

\section{Разность степеней и её знак}

Если $a>0$, $a\ne1$, то
\[
\operatorname{sgn}\bigl(a^u-a^v\bigr)
=\operatorname{sgn}\bigl((a-1)(u-v)\bigr).
\]
Это правило применимо именно к разности двух степеней одного основания. Дополнительное слагаемое или знак плюса внутри скобки разрушает эту структуру.

В предыдущем примере
\[
\operatorname{sgn}(25-5^x)=\operatorname{sgn}(2-x),
\qquad
\operatorname{sgn}(2^x-1)=\operatorname{sgn}(x).
\]
Следовательно,
\[
(25-5^x)(2^x-1)>0
\iff x(2-x)>0
\iff\boxed{x\in(0;2)}.
\]

\begin{center}
\begin{tikzpicture}[x=2.1cm,y=1cm]
  \draw[->,textgray] (-2.3,0)--(3.3,0);
  \fill[accent] (0,0) circle (2pt) node[below=5pt] {$0$};
  \fill[accent] (2,0) circle (2pt) node[below=5pt] {$2$};
  \node[above] at (-1.1,0) {$-$};
  \node[above] at (1,0) {$+$};
  \node[above] at (2.65,0) {$-$};
  \draw[accent,line width=2pt] (0.03,0.15)--(1.97,0.15);
\end{tikzpicture}
\end{center}

\section{Когда разрешено делить неравенство}

Деление на выражение с переменной допустимо только после установления его знака. Если делитель положителен, знак сохраняется; если отрицателен, меняется; если может обращаться в нуль, нужно рассматривать случаи.

В показательных задачах полезен факт
\[
a^x>0\qquad\text{для всех }x\in\R\text{ при }a>0.
\]
Поэтому на $C\,a^{f(x)}$ можно делить без смены знака, если $C>0$. Это обоснование следует написать до деления.

{\small\setstretch{1.15}
\section{Типичные ошибки}

\begin{enumerate}
  \item Снимать разные основания, не приведя их к одной базе.
  \item Не проверять, больше или меньше единицы основание.
  \item Записывать систему вместо совокупности при $|f(x)|\ge g(x)$.
  \item Считать, что отрицательный дискриминант автоматически означает отсутствие решений неравенства.
  \item Делить на выражение с переменной без проверки его знака и возможности нуля.
  \item Применять правило разности степеней к сумме или к выражению с лишним слагаемым.
  \item После группировки потерять минус перед второй скобкой.
  \item Включать граничные точки в ответ при строгом неравенстве.
\end{enumerate}

\section{Итоговый чек-лист}

\begin{itemize}[label=$\square$,itemsep=1pt]
  \item Я привёл сравниваемые степени к одинаковому основанию.
  \item Я установил: $a>1$ или $0<a<1$.
  \item Я корректно сохранил или изменил знак неравенства.
  \item Я изолировал модуль перед равносильным переходом.
  \item Я различаю систему и совокупность.
  \item При четырёх слагаемых я проверил возможность группировки.
  \item Я применяю правило разности степеней только к $a^u-a^v$.
  \item Перед делением я доказал знак делителя.
  \item Я проверил границы, ОДЗ и скобки в ответе.
\end{itemize}
}

\newpage
\section{Тренировочный блок}

\textbf{Средний уровень}
\begin{enumerate}
  \item $\left(\dfrac32\right)^{|x|-4}\ge1$.
  \item $\left(\dfrac23\right)^{x^2-5x+6}\ge1$.
  \item $|x+1|-3x+5\ge0$.
\end{enumerate}

\textbf{Уровень выше среднего}
\begin{enumerate}[resume]
  \item $|x+1|-3x+5<0$.
  \item $16\cdot3^x-12^x+4^x-16>0$.
  \item $9\cdot2^x-6^x+3^x-9\ge0$.
  \item $(5^x-25)(7^x-1)\le0$.
  \item $\dfrac{2^x-4}{3^x-1}>0$.
  \item $\left(\dfrac73\right)^{x-1}\ge1$.
\end{enumerate}

\textbf{Сложный уровень}
\begin{enumerate}[resume]
  \item $81\cdot2^x-18^x+9^x-81\le0$.
\end{enumerate}

\newpage
\section{Ответы}

\begin{table}[ht]
\centering
\caption{Ответы к тренировочному блоку}
\begin{tabularx}{\linewidth}{>{\bfseries}p{12mm} X}
\toprule
№ & Ответ \\
\midrule
1 & $(-\infty;-4]\cup[4;+\infty)$ \\
2 & $[2;3]$ \\
3 & $(-\infty;3]$ \\
4 & $(3;+\infty)$ \\
5 & $(0;2)$ \\
6 & $[0;2]$ \\
7 & $[0;2]$ \\
8 & $(-\infty;0)\cup(2;+\infty)$ \\
9 & $[1;+\infty)$ \\
10 & $(-\infty;0]\cup[2;+\infty)$ \\
\bottomrule
\end{tabularx}
\end{table}

\vfill
\begin{center}
\textcolor{softgray}{Сначала восстановите структуру выражения, затем выполняйте вычисления.}
\end{center}

\end{document}
